% ID: SEL-D0011-BA-E01
% ESTADO: cotejada. Segunda pasada visual completa y borrador compilado revisados.
% Fuente: Selectividad/matematicas_aplicadas/EXTRA_SUPLENTE_Modelo-1-Enunciados (Mat-Ap-2021).pdf
% SHA-256: 4a6c7d34ff35c3e7678c52970b4d4558cee2c35224b46c66035cb0bc8d5a7ec8
% Página: 1; bloque A; ejercicio original 1.

\textbf{(2.5 puntos)} Un laboratorio farmacéutico tiene una línea de producción
con dos medicamentos $A$ y $B$, con marca comercial y genérico respectivamente,
de los cuales, entre los dos como máximo puede fabricar 10 unidades a la hora.
Desde el punto de vista del rendimiento, se han de producir al menos 4 unidades
por hora entre los dos y por motivos de política sanitaria, la producción
de $A$ ha de ser como mucho 2 unidades más que la de $B$.

Cada unidad de tipo $A$ que vende le produce un beneficio de 60 euros,
mientras que cada unidad de tipo $B$ le produce un beneficio de 25 euros.
Si se vende todo lo que se produce, determine las unidades de cada medicamento
que deberá fabricar por hora para maximizar su beneficio y obtenga el valor
de dicho beneficio.
